\documentclass{article}

\usepackage{../handouts}

\title{CSCI 432 Handout 10: Topological Sort}
\author{Name: \rule{3in}{0.15mm}}

\begin{document}
\maketitle

\section*{Topological Sort: The Algorithm}

Topological sort takes a poset $(P,\leq)$ and returns a total order of all
objects that is \emph{compatible} with the partial order of the poset.  By
compatible, we mean that for all $a \leq b \in P$, we have $a$ appears before
$b$ in the total order.  For example, consider $P$ to be a set of points in
$\R^2$ and let the partial order be the product order; that is, $(a_x,a_y) \leq
(b_x,b_y)$ if and only if $a_x \leq b_x$ and $a_y \leq b_y$.  Given a set:
$$ P = \{ (0,0), (\pi, \pi), (2,4), (1,2), (3, 2) \}$$
we can return the following total order:
$$P = [ (0,0), (1,2), (2,4), (3,2), (\pi, \pi)].$$
Note that this is not the only compatible total order.

\begin{enumerate}
    \item What other total order
        could have been returned?

    \textbf{Answer}
    \vspace{1in}

    \item Another order we can impose on $\R^2$ is the lexicographical ordering (sometimes
        called the dictionary ordering), where
        $(a_x,a_y) \leq (b_x,b_y)$ if and only if $a_x < a_y$ or ($a_x=a_y$ and $a_y
        \leq b_y$).  What total orders on $P$ exist that are compatible with the
        lexicographical ordering?

    \textbf{Answer}
\end{enumerate}

\newpage
The following algorithm takes as input a poset, represented as a directed
acyclic graph, and returns a compatible total order.  If the poset orders $a
\leq b$, then the DAG has an edge from vertex $a$ to vertex $b$.  For this
algorithm, vertices have at least one attribute: outgoing, which stores a list
of vertices that define the outgoing edges.

\begin{algorithm}[h]
    \caption{Topological Sort}
    %
    \begin{algorithmic}[1]
        \REQUIRE $G=(V,E)$, a DAG

        \ENSURE $L$, a list of vertices in G in a total order compatible with
        the partial order of the DAG.

        \FOR {$v \in V$}
            \STATE Add an attribute `$v$.count' and initialize it to the indegree of $v$.
        \ENDFOR

        \STATE $Q \gets$ queue containing the vertices with `count'=0.\label{ln:initQ}

        \STATE $L \gets$ array of length $|V|$
        \STATE $i \gets 1$

        \WHILE {$Q$ is not empty}
            \STATE $v \gets Q.\textsc{POP}$
            \STATE $L[i] \gets v$
            \STATE $i++$
            \FOR {$u \in v$.outgoing}
                \STATE $u$.count$--$
                \IF {$u$.count $=0$}
                    \STATE Add $u$ to $Q$
                \ENDIF
            \ENDFOR
        \ENDWHILE

        \RETURN $L$
    \end{algorithmic}
\end{algorithm}

\pagebreak

\begin{enumerate}\setcounter{enumi}{2}

    \item In Line~\ref{ln:initQ}, why is $Q$ not empty? (To answer this fully, you'll
        need to write a proof).

        \textbf{Answer}
        \vspace{1in}

    \item   If you are familiar with dynamic programming:
        Explain how topological sort is helpful for dynamic programming.

        \textbf{Answer}
        \vspace{1in}

    \item Walk through a small example.

        \textbf{Answer}
\end{enumerate}

\pagebreak
\section*{Topological Sort: The Decrementing Function}

Recall from HO-07 that a decrementing function is a function $D \colon
\mathcal{S} \to \N$, where $\mathcal{S}$ is the state space that
such that (when there is no recursion) each time a loop is re-entered, the function is strictly less than
the last time it entered that loop.


\begin{enumerate}\setcounter{enumi}{5}
    \item What is the decrementing function for the for loop?

    \textbf{Answer}
    \vspace{1in}

    \item What is the decrementing function for the while loop?

    \textbf{Answer}
    \vspace{1in}
\end{enumerate}

\newpage
\section*{Topological Sort: The Loop Invariant}

Next, we investigate partial correctness of the while loop.  For this loop, what
are the following statements (start with your best guess, then come back and
revise if needed):

The loop guard $G$:
\vspace{1.5in}

The post-condition $Q$:
\vspace{1.5in}

The pre-condition $P$:
\vspace{1.5in}

The loop invariant $L=L_i$:
\vspace{1.5in}

\newpage
Can you use this loop invariant
to prove partial correctness.

\begin{enumerate}\setcounter{enumi}{7}
    \item Initialization: If the preconditions are met and if we enter the loop, then the
    loop invariant is correct. ($P \wedge G \implies L$).

    \textbf{Answer}
    \vspace{2in}


    \item Maintenance: If we entered the loop and the loop invariant held, then when we
    exit the loop, the loop invariant will still hold: ($G \wedge L_i \implies
    L_{i+1})$.

    \textbf{Answer}
    \vspace{2in}


    \item End: If the loop ends and the loop invariant was correct, then the algorithm is
    correct. ($L \wedge \neg G \implies Q$).


    \textbf{Answer}
    \vspace{2in}
\end{enumerate}

\end{document}
